% =====================================================================
%  estilo.tex — Preámbulo y estilo visual de la memoria
%  No hace falta tocar este archivo para escribir la memoria.
%
%  Sistema visual:
%    - Texto en Palatino (serifa), todo lo "estructural" (títulos,
%      etiquetas, pies, cabeceras, índice) en Source Sans, y el código en
%      Source Code Pro.
%    - Tinta casi negra, una escala de grises y un único color de acento,
%      reservado para números, etiquetas y filetes clave.
% =====================================================================

% --- Idioma y codificación -------------------------------------------
\usepackage[T1]{fontenc}
\usepackage[utf8]{inputenc}
\usepackage[spanish,es-tabla,es-nodecimaldot,es-noshorthands]{babel}

% --- Tipografía ------------------------------------------------------
\usepackage{newpxtext,newpxmath}
\usepackage[semibold,scale=0.95]{sourcesans}
\usepackage[scale=0.84]{sourcecodepro}
\usepackage[protrusion=true, expansion=true, kerning=true]{microtype}
\linespread{1.12}

% Pesos de Source Sans
\newcommand{\sansligera}{\sffamily\fontseries{l}\selectfont}
\newcommand{\sansextraligera}{\sffamily\fontseries{el}\selectfont}
\newcommand{\sansnormal}{\sffamily\fontseries{m}\selectfont}
\newcommand{\sansseminegra}{\sffamily\fontseries{sb}\selectfont}

% Etiqueta: mayúsculas espaciadas en Source Sans seminegra.
%   \etiqueta{texto}   (el tamaño lo pone el contexto)
\newcommand{\etiqueta}[1]{%
  {\sansseminegra\textls[110]{\MakeUppercase{#1}}}}

% --- Página ----------------------------------------------------------
\usepackage[a4paper, top=3cm, bottom=3cm, inner=3.2cm, outer=3.2cm,
            headheight=16pt, headsep=26pt, footskip=36pt]{geometry}
\setlength{\parskip}{0.55em plus 0.1em}
\setlength{\parindent}{0pt}
\raggedbottom
% Sin líneas viudas ni huérfanas.
\widowpenalty=10000
\clubpenalty=10000
\displaywidowpenalty=10000

% --- Colores ----------------------------------------------------------
\usepackage[table]{xcolor}
\definecolor{principal}{HTML}{1B2129}   % tinta
\definecolor{acento}{HTML}{1F5A8B}      % azul acero, único color
\definecolor{acentosuave}{HTML}{E8EFF6}
\definecolor{gristexto}{HTML}{545B64}
\definecolor{grissuave}{HTML}{8A9099}
\definecolor{grislinea}{HTML}{CDD1D6}
\definecolor{grisfondo}{HTML}{F7F8F9}
\definecolor{grisbarra}{HTML}{EEF0F2}
\definecolor{codclave}{HTML}{1F5A8B}
\definecolor{codcadena}{HTML}{4D6B3C}
\definecolor{codcoment}{HTML}{8A9099}

% --- Paquetes generales ----------------------------------------------
\usepackage{amsmath}
\usepackage{graphicx}
\graphicspath{{figuras/}}
\usepackage{booktabs,tabularx,array,multirow}
\setlength{\heavyrulewidth}{0.07em}
\setlength{\lightrulewidth}{0.035em}
\setlength{\cmidrulewidth}{0.035em}
\renewcommand{\arraystretch}{1.18}
\usepackage{float}
% Cabeceras de las tablas (\textbf) en Source Sans seminegra.
\newcommand{\textcabecera}[1]{{\sansseminegra #1}}
\AtBeginEnvironment{table}{\let\textbf\textcabecera}
\usepackage{needspace}
% Pies: "FIGURA 3" en etiqueta de acento y el texto en gris.
\usepackage{caption}
\DeclareCaptionLabelFormat{etiquetada}{\etiqueta{#1~#2}}
\DeclareCaptionLabelSeparator{espacio}{\hspace{0.8em}}
\captionsetup{font={small,color=gristexto}, labelfont={footnotesize,color=acento},
              labelformat=etiquetada, labelsep=espacio, skip=9pt,
              margin=0.6cm, justification=justified, singlelinecheck=true}
\captionsetup[table]{position=top, skip=7pt}
\usepackage{subcaption}
\usepackage{enumitem}
\setlist{itemsep=0.3em, topsep=0.4em, parsep=0pt, leftmargin=1.7em}
\setlist[itemize,1]{label={\color{acento}\rule[0.27em]{0.45em}{0.6pt}}}
\setlist[itemize,2]{label={\color{grissuave}\textperiodcentered}}
\setlist[enumerate,1]{label={\color{acento}\sansseminegra\arabic*}}
\usepackage{tikz}
\usepackage[most]{tcolorbox}
\usepackage{listings}
\usepackage{fancyhdr}
\usepackage{titlesec}
\usepackage[titles]{tocloft}
\usepackage{etoolbox}
\usepackage{bookmark}

% --- Títulos ---------------------------------------------------------
% Ejercicio / tarea / apartado: etiqueta de acento, título en Source Sans
% y filete fino a todo lo ancho.
\titleformat{\section}[display]
  {\normalfont\color{principal}}
  {\footnotesize\color{acento}\etiqueta{\NombreSeccion~\thesection}}
  {0.55em}
  {\sansnormal\fontsize{21}{25}\selectfont}
  [\vspace{0.75em}{\color{grislinea}\titlerule[0.4pt]}]
% Secciones sin número (Introducción, Planteamiento, Conclusiones).
\titleformat{name=\section, numberless}[display]
  {\normalfont\color{principal}}
  {}
  {0em}
  {\sansnormal\fontsize{21}{25}\selectfont}
  [\vspace{0.75em}{\color{grislinea}\titlerule[0.4pt]}]
\titleformat{\subsection}
  {\normalfont\sansseminegra\large\color{principal}}
  {{\sansnormal\color{acento}\thesubsection}}{0.85em}{}
\titleformat{\subsubsection}
  {\normalfont\sansseminegra\normalsize\color{gristexto}}
  {}{0em}{}
\titlespacing*{\section}{0pt}{2.4em}{1.3em}
\titlespacing*{\subsection}{0pt}{1.7em}{0.5em}
\titlespacing*{\subsubsection}{0pt}{1.2em}{0.3em}

% Nombre de las secciones numeradas en cada práctica.
\newcommand{\NombreSeccion}{Ejercicio}
\newcommand{\NombreToc}{Ejercicio}
\newcommand{\PracticaActual}{\Titulo}

% Anclas de hyperref únicas aunque la numeración se reinicie en cada práctica.
\renewcommand{\theHsection}{\arabic{part}.\arabic{section}}
\renewcommand{\theHsubsection}{\theHsection.\arabic{subsection}}
\renewcommand{\theHsubsubsection}{\theHsubsection.\arabic{subsubsection}}

% --- Índice ----------------------------------------------------------
\renewcommand{\contentsname}{Índice}
\renewcommand{\cfttoctitlefont}{\normalfont\sansnormal\fontsize{21}{25}\selectfont\color{principal}}
\renewcommand{\cftaftertoctitle}{\par\vspace{0.75em}{\color{grislinea}\rule{\linewidth}{0.4pt}}}
\setlength{\cftbeforetoctitleskip}{0pt}
\setlength{\cftaftertoctitleskip}{1.6em}
% Prácticas (la entrada completa se escribe en \practica)
\renewcommand{\cftpartfont}{\normalfont\sansseminegra\color{principal}}
\renewcommand{\cftpartpagefont}{\normalfont\sansseminegra\color{principal}}
\setlength{\cftbeforepartskip}{2.1em}
% Ejercicios, tareas y apartados
\renewcommand{\cftsecfont}{\normalfont\color{principal}}
\renewcommand{\cftsecpagefont}{\normalfont\color{principal}}
\renewcommand{\cftsecpresnum}{\begingroup\footnotesize\color{grissuave}%
  \sansseminegra\lsstyle\MakeUppercase{\NombreToc}~}
\renewcommand{\cftsecaftersnum}{\endgroup}
\setlength{\cftsecnumwidth}{7.6em}
\setlength{\cftbeforesecskip}{0.5em}
\renewcommand{\cftsecleader}{{\color{grislinea}\cftdotfill{\cftdotsep}}}
\renewcommand{\cftdotsep}{2.6}
% Apartados internos
\renewcommand{\cftsubsecfont}{\normalfont\color{gristexto}}
\renewcommand{\cftsubsecpagefont}{\normalfont\color{gristexto}}
\setlength{\cftsubsecindent}{7.6em}
\setlength{\cftsubsecnumwidth}{2.4em}
\setcounter{tocdepth}{1}

% --- Cabecera y pie ----------------------------------------------------
% Cabecera sin filete: práctica y ejercicio a la izquierda, página a la
% derecha. El pie queda vacío.
\pagestyle{fancy}
\fancyhf{}
\renewcommand{\headrulewidth}{0pt}
\fancyhead[L]{\footnotesize\color{acento}\etiqueta{\PracticaActual}%
  \hspace{0.9em}{\color{grislinea}\rule[-0.15em]{0.4pt}{0.95em}}\hspace{0.9em}%
  {\sansnormal\color{grissuave}\nouppercase{\rightmark}}}
\fancyhead[R]{\footnotesize\sansseminegra\color{principal}\thepage}
\renewcommand{\sectionmark}[1]{\markright{#1}}
\renewcommand{\subsectionmark}[1]{}
\fancypagestyle{plain}{\fancyhf{}\renewcommand{\headrulewidth}{0pt}%
  \fancyfoot[R]{\footnotesize\sansseminegra\color{grissuave}\thepage}}

% --- Código Python -------------------------------------------------------
\lstdefinestyle{python}{
  language=Python,
  basicstyle=\ttfamily\small\linespread{1.08}\selectfont\color{principal},
  keywordstyle=\color{codclave}\bfseries,
  stringstyle=\color{codcadena},
  commentstyle=\color{codcoment}\itshape,
  numberstyle=\scriptsize\ttfamily\color{grissuave},
  numbers=left, numbersep=13pt,
  showstringspaces=false, breaklines=true, tabsize=4,
  columns=fullflexible, keepspaces=true,
  morekeywords={as,True,False,None,with},
  literate=
    {á}{{\'a}}1 {é}{{\'e}}1 {í}{{\'i}}1 {ó}{{\'o}}1 {ú}{{\'u}}1
    {Á}{{\'A}}1 {É}{{\'E}}1 {Í}{{\'I}}1 {Ó}{{\'O}}1 {Ú}{{\'U}}1
    {ñ}{{\~n}}1 {Ñ}{{\~N}}1 {ü}{{\"u}}1 {¿}{{?`}}1 {¡}{{!`}}1
    {€}{{\texteuro}}1 {·}{{$\cdot$}}1 {σ}{{$\sigma$}}1,
}
\lstset{style=python}
\renewcommand{\lstlistingname}{Código}

% Bloques de código: fondo gris muy claro, barra superior con el nombre
% del archivo y filete fino abajo.
\lstdefinestyle{bloque}{
  style=python,
  backgroundcolor=\color{grisfondo},
  frame=b, framerule=0.4pt, rulecolor=\color{grislinea},
  xleftmargin=3em, framexleftmargin=3em,
  xrightmargin=0pt, framexrightmargin=0pt,
  framextopmargin=8pt, framexbottommargin=8pt,
  aboveskip=1.3em, belowskip=1.3em, captionpos=t,
}
\DeclareCaptionFormat{barracodigo}{%
  {\setlength{\fboxsep}{0pt}\colorbox{grisbarra}{%
    \hbox to \linewidth{\rule[-0.6em]{0pt}{1.75em}\hspace{11pt}%
      {\color{acento}\rule[-0.1em]{2pt}{0.85em}}\hspace{8pt}%
      {\ttfamily\footnotesize\color{principal}#3}\hfil
      {\scriptsize\color{grissuave}\etiqueta{Python}}\hspace{11pt}}}}}
\captionsetup[lstlisting]{format=barracodigo, skip=0pt, margin=0pt,
  singlelinecheck=false}

% Bloque de código en el texto:
%   \begin{python}{archivo.py} ... \end{python}
%   (con {} vacío no se muestra la barra del archivo)
\lstnewenvironment{python}[1]
  {\endgraf\Needspace{10\baselineskip}%
   \ifstrempty{#1}%
     {\lstset{style=bloque, frame=tb}}%
     {\lstset{style=bloque, title={#1}}}}
  {}

% Código completo desde un archivo de la carpeta codigo/:
%   \codigo{archivo.py}
\newcommand{\codigo}[1]{%
  \lstinputlisting[style=bloque, title={#1}]{codigo/#1}}

% Los lstlisting sueltos tampoco se parten si son cortos.
\BeforeBeginEnvironment{lstlisting}{\par\Needspace{12\baselineskip}}

% --- Cajas -------------------------------------------------------------

% Enunciado: fondo de acento muy suave y filete de acento a la izquierda.
\newtcolorbox{enunciado}[1][]{
  enhanced, frame hidden, colback=acentosuave!55!white,
  borderline west={1.2pt}{0pt}{acento}, sharp corners,
  left=14pt, right=14pt, top=10pt, bottom=10pt, boxsep=0pt,
  fontupper=\itshape\linespread{1.1}\selectfont, coltext=gristexto,
  before upper={{\upshape\footnotesize\color{acento}\etiqueta{Enunciado}}\par
                \vspace{0.35em}},
  before skip=1em, after skip=1.4em, #1
}

% Conclusión: recuadro de línea fina con etiqueta.
\newtcolorbox{conclusion}[1][]{
  enhanced, breakable, colback=white, colframe=acento,
  boxrule=0.5pt, sharp corners,
  left=14pt, right=14pt, top=10pt, bottom=10pt,
  before upper={{\footnotesize\color{acento}\etiqueta{Conclusión}}\par
                \vspace{0.35em}},
  before skip=1.3em, after skip=1.3em, #1
}

% Nota o aclaración: texto pequeño con filete gris.
\newtcolorbox{nota}[1][]{
  enhanced, breakable, frame hidden, interior hidden,
  borderline west={0.8pt}{0pt}{grislinea},
  left=14pt, right=0pt, top=1pt, bottom=1pt, boxsep=0pt,
  fontupper=\small, coltext=gristexto,
  before upper={{\scriptsize\color{grissuave}\etiqueta{Nota}}\hspace{0.8em}},
  before skip=0.7em, after skip=1.2em, #1
}

% --- Comandos de ayuda -------------------------------------------------

% Texto pendiente de escribir (se ve en gris; bórralo al redactar).
\newcommand{\pendiente}[1]{%
  {\color{grissuave}\itshape\small[\,#1\,]}}

% Hueco para una figura que todavía no existe (compila sin la imagen).
%   \figurapendiente{Descripción de la figura}
\newcommand{\figurapendiente}[1]{%
  \begin{figure}[H]\centering
  \tikz\node[draw=grislinea, dashed, line width=0.5pt, fill=grisfondo,
             minimum width=0.82\linewidth, minimum height=4.2cm,
             text width=0.7\linewidth, align=center,
             font=\small\color{grissuave}]
             {\hyphenpenalty=10000 {\scriptsize\etiqueta{Figura pendiente}}\\[0.5em]
              \itshape #1};
  \end{figure}}

% Figura: \figura[ancho]{archivo}{pie}{etiqueta}
\newcommand{\figura}[4][0.9\linewidth]{%
  \begin{figure}[H]\centering
  \includegraphics[width=#1]{#2}
  \caption{#3}\label{#4}
  \end{figure}}

% Nombre de variable o función en línea: \var{EstimatedSalary}
\newcommand{\var}[1]{\texttt{#1}}

% Cabecera de ejercicio, tarea o apartado:
%   \ejercicio{1}{Título}{scipy-1.py}
%   \ejercicio[opcional]{4}{Título}{practica2.py}
%   \ejercicio{7}{Título}{}            (sin archivo)
% Pregunta del enunciado al principio de cada respuesta:
%   \item \pregunta{¿Qué escalado elige?} Texto de la respuesta...
\newcommand{\pregunta}[1]{{\sansseminegra\color{principal}#1}\par\nopagebreak\vspace{0.2em}}

\newif\ifsinsalto
\newcommand{\ejercicio}[4][]{%
  \ifsinsalto\global\sinsaltofalse\else\clearpage\fi
  \setcounter{section}{\numexpr#2-1\relax}%
  \section{#3}
  \ifboolexpr{test {\ifstrempty{#4}} and test {\ifstrempty{#1}}}{}{%
    \vspace{-0.55em}
    {\footnotesize\color{grissuave}%
     \ifstrempty{#4}{}{\etiqueta{Archivo}\hspace{0.7em}%
       {\color{gristexto}\texttt{#4}}}%
     \ifboolexpr{test {\ifstrempty{#4}} or test {\ifstrempty{#1}}}{}%
       {\hspace{1.4em}}%
     \ifstrempty{#1}{}{{\setlength{\fboxsep}{2.2pt}%
       \colorbox{acentosuave}{\scriptsize\color{acento}\etiqueta{#1}}}}}%
    \par\vspace{0.4em}}}

% Portadilla de cada práctica (página propia):
%   \practica{sesión}{número}{título}{descripción}{material}{Ejercicio}
\newcommand{\practica}[6]{%
  \clearpage
  \setcounter{part}{\numexpr#1-1\relax}%
  \refstepcounter{part}%
  \setcounter{section}{0}%
  \gdef\PracticaActual{Práctica~#2}%
  \gdef\NombreSeccion{#6}%
  \addtocontents{toc}{\protect\gdef\protect\NombreToc{#6}}%
  \phantomsection
  \addcontentsline{toc}{part}{\protect\makebox[7.6em][l]%
    {\protect\footnotesize\protect\color{acento}\protect\etiqueta{Práctica~#2}}%
    #3}%
  \markright{}%
  \thispagestyle{empty}%
  \begingroup\raggedright\color{principal}
  \vspace*{0.1\textheight}
  {\footnotesize\color{acento}\etiqueta{Práctica~#2}\par}
  \vspace{0.2em}
  {\sansextraligera\fontsize{150}{150}\selectfont\color{acento}\hspace*{-0.06em}#2\par}
  \vspace{1.2em}
  {\sansnormal\fontsize{30}{35}\selectfont #3\par}
  \vspace{1.6em}
  {\color{acento}\rule{2.2cm}{1.4pt}\par}
  \vspace{1.6em}
  \begin{minipage}{0.78\linewidth}
    \raggedright\linespread{1.18}\selectfont\color{gristexto}#4
  \end{minipage}\par
  \vfill
  {\color{grislinea}\rule{\linewidth}{0.4pt}\par}
  \vspace{0.9em}
  {\scriptsize\color{grissuave}\etiqueta{Material}\par}
  \vspace{0.45em}
  {\small\color{gristexto}#5\par}
  \endgroup
  \clearpage}

% Separador de bloque dentro de una práctica (por ejemplo, por librería):
%   \bloque{SciPy}{Ejercicios 1 a 5}
\newcommand{\bloquetoc}[1]{%
  \par\addvspace{1em}%
  {\scriptsize\color{acento}\etiqueta{#1}\par}\nopagebreak[4]%
  \vspace{0.15em}}
\newcommand{\bloque}[2]{%
  \clearpage\global\sinsaltotrue
  \addtocontents{toc}{\protect\bloquetoc{#1}}%
  \noindent{\footnotesize\color{acento}\etiqueta{#1}}\hspace{0.9em}%
  {\color{grislinea}\leaders\vrule height 0.33em depth -0.29em\hfill}%
  \hspace{0.9em}{\footnotesize\sansnormal\color{grissuave}#2}\par
  \nopagebreak\vspace{-0.6em}}

% --- Enlaces (siempre al final del preámbulo) ------------------------
\usepackage{hyperref}
\hypersetup{colorlinks=true, linkcolor=principal, urlcolor=acento,
            citecolor=principal, bookmarksnumbered=true,
            bookmarksopen=true, pdfdisplaydoctitle=true,
            pdfauthor={\Autor}, pdftitle={\Titulo\ de laboratorio. \Subtitulo},
            pdfsubject={\Asignatura}}
\usepackage[spanish,noabbrev]{cleveref}
\crefname{figure}{figura}{figuras}
\crefname{table}{tabla}{tablas}
\crefname{lstlisting}{código}{códigos}
\crefname{listing}{código}{códigos}
\Crefname{listing}{Código}{Códigos}
\Crefname{lstlisting}{Código}{Códigos}
